\subsection{From linear combinations to a plane}

In the vector chapter we studied all combinations
\[
s\mathbf u+t\mathbf v.
\]
When $\mathbf u$ and $\mathbf v$ are not parallel, these vectors form a
plane through the origin. Starting from a fixed point $P_0$ gives
\[
\Pi=\{P_0+s\mathbf u+t\mathbf v:s,t\in\mathbb R\}.
\]
Thus a plane is built from the same two operations already known for vectors:
scaling and addition.

The previous chapter proved that
\[
(\mathbf u\times\mathbf v)\cdot(a\mathbf u+b\mathbf v)=0
\]
for all real $a,b$. Therefore
\[
\mathbf n=\mathbf u\times\mathbf v
\]
is perpendicular to every displacement generated by $\mathbf u$ and
$\mathbf v$. This is why the cross product becomes a normal vector.

\subsection{The normal equation}

A point $P$ belongs to the plane through $P_0$ exactly when $P-P_0$ is an
allowed displacement in the plane. Hence
\[
\boxed{\mathbf n\cdot(P-P_0)=0.}
\]
For
\[
P=(x,y,z),\qquad P_0=(x_0,y_0,z_0),\qquad \mathbf n=(a,b,c),
\]
this becomes
\[
a(x-x_0)+b(y-y_0)+c(z-z_0)=0,
\]
or
\[
ax+by+cz=d.
\]
The coefficients $(a,b,c)$ are the coordinates of the normal vector.

\begin{center}
\begin{tikzpicture}[scale=0.9,>=stealth]
    \coordinate (A) at (0.4,1.0);
    \coordinate (B) at (4.9,1.0);
    \coordinate (C) at (6.0,3.4);
    \coordinate (D) at (1.5,3.4);
    \coordinate (P0) at (2.2,2.0);
    \coordinate (P) at (4.3,2.8);
    \coordinate (N) at (2.2,5.0);
    \filldraw[fill=geometricSubsetColor!10,draw=geometricSubsetColor,thick]
        (A)--(B)--(C)--(D)--cycle;
    \draw[blue!70!black,line width=1.1pt,->] (P0)--(P)
        node[midway,below right] {$P-P_0$};
    \draw[orange!85!black,line width=1.1pt,->] (P0)--(N)
        node[midway,left] {$\mathbf n$};
    \fill[geometricSubsetColor] (P0) circle (2.2pt) node[below left] {$P_0$};
    \fill[geometricSubsetColor] (P) circle (2.2pt) node[above right] {$P$};
    \node[subset label] at (5.2,3.0) {$\Pi$};
\end{tikzpicture}
\end{center}

\subsection{Distance from a point to a plane}

Only the component of $P-P_0$ in the normal direction contributes to the
shortest distance. Its signed scalar component is
\[
\frac{\mathbf n\cdot(P-P_0)}{\|\mathbf n\|}.
\]
Therefore
\[
\boxed{
d(P,\Pi)=\frac{|\mathbf n\cdot(P-P_0)|}{\|\mathbf n\|}.
}
\]
For $ax+by+cz=d$ and $P=(x_P,y_P,z_P)$,
\[
\boxed{
d(P,\Pi)=
\frac{|ax_P+by_P+cz_P-d|}{\sqrt{a^2+b^2+c^2}}.
}
\]
This is the length of an orthogonal projection, not an isolated formula.
